\subsection{Вес и сила реакции  опоры}
%%%%%%%%%%%%%%%%%%%%%%%%%%%%%%%%%%%%%%%%%%%%%%%%%%%%
%%%%%%%%%%%%%%%%%%%%%%%%%%%%%%%%%%%%%%%%%%%%%%%%%%%%

 \z[{\begin{wrapfigure}[5]{r}{0.15\textwidth}%
		\vspace{-15pt}%
		\begin{flushright}%
 			\includegraphics[width=0.1\textwidth]{4.4.1}%
		\end{flushright}%
	\end{wrapfigure}%
	}]%
	{Маляр работает в подвесной люльке. Ему понадобилось срочно подняться вверх. Он принимается тянуть за веревку с такой силой, что сила его давления на пол люльки уменьшилась до $400$ H. Масса люльки $12$ кг, масса маляра $72$ кг. Чему равно ускорение люльки?}
 %
{$a=3.5\,\text{м}/\text{с}^2$}
%%%%%%%%%%%%%%%%%%%%%%%%%%%%%%%%%%%%%%%%%%%%%%%%%%%% 
 \z[{\begin{wrapfigure}[4]{r}{0.25\textwidth}%
		\vspace{-15pt}%
		\begin{flushright}%
 			\includegraphics[width=0.25\textwidth]{4.4.2}%
		\end{flushright}%
	\end{wrapfigure}%
	}]%
	{Сила $F$ приводит в движение клин 1 и штифт 2. Угол наклона клина равен $\alpha$, масса клина равна $m$, масса штифта тоже равна $m$, трение отсутствует. Найти ускорение и силу взаимодействия клина и штифта.}
 %
{$a=\frac{(F\cos\alpha-mg\sin\alpha)\cos\alpha}{m};$~ $F_\text{взаим} = F\sin\alpha+mg\cos\alpha$}
%%%%%%%%%%%%%%%%%%%%%%%%%%%%%%%%%%%%%%%%%%%%%%%%%%%% 
 \z[{\begin{wrapfigure}[5]{r}{0.25\textwidth}%
		\vspace{+10pt}%
		\begin{flushright}%
 			\includegraphics[width=0.25\textwidth]{4.4.3}%
		\end{flushright}%
	\end{wrapfigure}%
	}]%
	{На гладкой горизонтальной поверхности лежит клин с углом $\alpha = 15\degree$ при основании, упираясь торцом в неподвижную вертикальную стенку. По верхней грани клина соскальзывает без трения тело массы $m = 0.20\,\text{кг}$ (см. рисунок). Найдите силу нормального давления $N$ клина на стенку.}
%
{$N=\frac{1}{2}mg\sin{2\alpha}=0.5\,\text{Н}$}
%%%%%%%%%%%%%%%%%%%%%%%%%%%%%%%%%%%%%%%%%%%%%%%%%%%%
 \z[{\begin{wrapfigure}[5]{r}{0.25\textwidth}%
		\vspace{-10pt}%
		\begin{flushright}%
 			\includegraphics[width=0.25\textwidth]{4.4.4}%
		\end{flushright}%
	\end{wrapfigure}%
	}]%
	{По вогнутому мосту, образующему дугу окружности радиусом $R$, движется вагонетка массой $M$. К вагонетке привязан трос длиной $L$, на конце которого закреплён груз массой $m$ (см. рисунок). В момент, когда вагонетка проходила нижнюю точку моста, трос был расположен вертикально, а скорости вагонетки и груза были равны $v_1$ и $v_2$ соответственно. Найдите в этот момент силу натяжения троса $T$ и силу $N$, с которой вагонетка давит на рельсы. Трос невесом и нерастяжим, трение не учитывать, размерами вагонетки и груза пренебречь.}
 %
{$T=m\left(g+\frac{v_1^2}{R}+\frac{(v_1-v_2)^2}{L}\right);$~ $N=(m+M)\left(g+\frac{v_1^2}{R}\right)+m\frac{(v_1-v_2)^2}{L}$}
%%%%%%%%%%%%%%%%%%%%%%%%%%%%%%%%%%%%%%%%%%%%%%%%%%%% 
 \z[{\begin{wrapfigure}[3]{r}{0.25\textwidth}%
		\vspace{-15pt}%
		\begin{flushright}%
 			\includegraphics[width=0.25\textwidth]{4.4.5}%
		\end{flushright}%
	\end{wrapfigure}%
	}]%
	{Между двумя одинаковыми гладкими брусками массы $m_1$ каждый вставлен клин массы $m_2$ с углом $\alpha$. Определите ускорение тел. }
%
{$a_1=\frac{m_2g\tg{\frac{\alpha}{2}}}{m_2+2m_1\tg^2{\frac{\alpha}{2}}};$~$a_2=\frac{m_2g}{m_2+2m_1\tg^2{\frac{\alpha}{2}}}$}
%%%%%%%%%%%%%%%%%%%%%%%%%%%%%%%%%%%%%%%%%%%%%%%%%%%%